\section{The first representation calculations}\label{sec:representations}

Write $\VV_\lambda$ for the irreducible complex representation
of $\mathfrak{so}(25)$ with highest weight $\lambda$ in the
orthogonal coordinates of $B_{12}$. Cohomology will sometimes
be centred by subtracting the complex dimension from its degree.
In this convention $h$ is the grading operator.

\subsection{The ten-dimensional comparison}
For $v_0^2=2$ and a $2v_0$-generic polarisation, the singular
moduli space has the O'Grady symplectic resolution
$\widetilde M$. This case is outside \eqref{eq:family} and
requires none of our domain analysis.

\begin{proposition}\label{prop:og10}
In the ten-dimensional case,
\[
 \IH^\bullet(M;\CC)=
 \VV_{(5)}\oplus\VV_{(3)}\oplus\VV_{(2,1)}\oplus\VV_{(1)}.
\]
\end{proposition}
\begin{proof}
The decomposition theorem calculated in \cite{dcrs21} gives
\[
 H^\bullet(\widetilde M)=\IH^\bullet(M)
 \oplus H^\bullet(\Sym^2M_0)[-2]
 \oplus H^\bullet(M_0)[-6],
\]
compatibly with the $H^\bullet(M)$ action. Each summand has
grading symmetric about the middle and is preserved by the
raising operators. It is therefore preserved by their lowering
operators \cite[Lemma~3.4]{fsy22}.

The diagonal restriction $H^2(M)\to H^2(M_0)$ is a
similitude. It is equivariant for locally trivial monodromy,
so its pullback of the Beauville--Bogomolov form is a scalar
multiple of $q_M$ by density. This multiple is nonzero on an
ample class. Since both spaces have dimension $23$, restriction
is an isomorphism. Thus the last summand is the LLV module
$\VV_{(2)}$ of $K3^{[2]}$, and the preceding summand is
$\Sym^2\VV_{(2)}$, with the indicated degree shifts.

By \cite{gklr22}, the $\mathfrak{so}(26)$ module of
$\widetilde M$ is $\VV_{(5)}\oplus\VV_{(2,2)}$.
Restriction to $\mathfrak{so}(25)$, the stabiliser of the
exceptional class, is
\[
 \VV_{(5)}\oplus\VV_{(4)}\oplus\VV_{(3)}\oplus2\VV_{(2)}
 \oplus\VV_{(1)}\oplus\VV_{(0)}\oplus\VV_{(2,2)}\oplus\VV_{(2,1)}.
\]
Subtract
$\Sym^2\VV_{(2)}=\VV_{(4)}\oplus\VV_{(2,2)}\oplus
\VV_{(2)}\oplus\VV_{(0)}$ and the last $\VV_{(2)}$.
Complete reducibility gives the stated decomposition.
\end{proof}

\subsection{A character in dimension eighteen}
Return to a Picard-rank-one quartic $S$, $v_0=(5,A,0)$,
and $M=M_A(2v_0)$. Then $M_0$ is of type $K3^{[3]}$.
Centre all gradings. Transport the LLV modules of Hilbert
schemes to $\mathfrak{so}(25)$ by complex similitudes of
their second cohomologies; their characters are independent
of these choices.

\begin{lemma}\label{lem:character}
The virtual character
$\operatorname{ch}H^\bullet(K3^{[9]})-
\operatorname{ch}\bigwedge^2H^\bullet(K3^{[3]})$
is the character of the module $\mathcal R$ whose nonzero
multiplicities are listed in Table \ref{tab:representation}.
Its dimension is $138065600$.
\end{lemma}

\begin{table}[ht]
\centering
\caption{The irreducible summands of $\mathcal R$. Omitted
coordinates of each highest weight are zero.}
\label{tab:representation}
\begin{tabular}{crcrcr}
\toprule
$\lambda$&multiplicity&$\lambda$&multiplicity&$\lambda$&multiplicity\\
\midrule
$(9)$&1&$(4,2)$&1&$(2,2)$&1\\
$(7,1)$&1&$(4,1,1)$&1&$(2,1,1)$&1\\
$(7)$&1&$(4,1)$&2&$(2,1)$&3\\
$(6,1)$&1&$(4)$&2&$(2)$&3\\
$(6)$&1&$(3,2)$&2&$(1,1,1)$&1\\
$(5,2)$&1&$(3,1)$&4&$(1,1)$&2\\
$(5,1)$&1&$(3)$&4&$(1)$&4\\
$(5)$&3&$(2,2,1)$&1&&\\
\bottomrule
\end{tabular}
\end{table}

\begin{proof}
Let $z_1,\ldots,z_{12}$ be the standard torus coordinates
for type $B_{12}$. The Hilbert-scheme character formula
of \cite[Theorem~1.2(1)]{gklr22} is
\begin{equation}\label{eq:Hilbert-character}
 \begin{split}
 1+q\sum_{i=1}^{12}(z_i+z_i^{-1})
 +\sum_{n\ge2}\operatorname{ch}H^\bullet(K3^{[n]})q^n
 =\prod_{m\ge1}\prod_{i=1}^{12}
 \frac{1}{(1-q^mz_i)(1-q^mz_i^{-1})}.
 \end{split}
\end{equation}
The convention here centres degrees. The identity is in the
formal power series ring over the $B_{12}$ representation ring;
its coefficient of $q$ is written separately because $K3^{[1]}$
does not have the same LLV algebra as $K3^{[n]}$ for $n\ge2$.
Expand through $q^9$, and use
\[
 \operatorname{ch}\bigwedge^2 V(z)
 =\tfrac12\big((\operatorname{ch}V(z))^2-
                    \operatorname{ch}V(z^2)\big).
\]
Successive subtraction of highest-weight characters gives the
multiplicities in the table. These are finite weight computations:
the positive roots are $e_i\pm e_j$ for $i<j$ and $e_i$,
and Freudenthal's recurrence determines each character from
its highest weight. The accompanying code performs this
calculation with exact integers. At $z=1$, the result is
\[
 143184000-\binom{3200}{2}=138065600,
\]
also the sum of the Weyl dimensions in the table. Replacing
the exterior square by the symmetric square gives negative
irreducible multiplicities. Repeating the calculation with
$K3^{[5]}$ and $K3^{[2]}$ gives Proposition \ref{prop:og10}.
\end{proof}

The character calculation is independent of whether it describes
$\IH(M)$. To state the geometric inputs separately, let
$\mathrm{BPS}_v$ be the BPS cohomology of one-dimensional
semistable sheaves of vector $v$ on a K3 surface, with the
normalisations of \cite{dhsm22}. We use the following named
hypotheses when indicated:
\begin{enumerate}
\item[(H1)] On a positive primitive ray $\ZZ_{>0}w$, the BPS
Lie algebra is freely generated by the intersection cohomologies
of the moduli spaces on that ray.
\item[(H2)] The graded BPS cohomology is independent of the
Euler characteristic when the first Chern class is fixed.
\end{enumerate}
These are the BPS and $\chi$-independence statements of the
preprints \cite{dhsm22,dhksv26,gwz26}. Under them, the degree
$2w$ part of the free Lie algebra and deformation equivalence
\cite[Theorem~1.7]{pr23} give a graded-vector-space isomorphism
\begin{equation}\label{eq:graded-comparison}
 H^\bullet(K3^{[9]})\cong\IH^\bullet(M)
                \oplus\bigwedge^2H^\bullet(K3^{[3]}).
\end{equation}
Indeed the primitive-vector moduli space on the ray has type
$K3^{[3]}$; changing the Euler characteristic of the doubled
vector to an odd integer makes it primitive and gives type
$K3^{[9]}$. The primitive cohomology is even, so its Lie bracket
contribution is the ordinary exterior square.

Consequently (H1)--(H2) imply vanishing of odd intersection
cohomology and the following Betti numbers through the middle:
\begin{equation}\label{eq:Betti}
\begin{array}{c|rrrrr}
 j&0&1&2&3&4\\\hline
 \dim\IH^{2j}&1&23&300&2876&22427\\[2pt]
 j&5&6&7&8&9\\\hline
 \dim\IH^{2j}&150053&887710&4727859&22432662&81617778
\end{array}
\end{equation}
The other numbers follow from
$\dim\IH^{36-2j}=\dim\IH^{2j}$. This is still a statement
about graded dimensions, not the LLV decomposition.

\subsection{Recovering a representation from a restricted character}
Choose coordinates $e_1,\ldots,e_{12}$ on a Cartan subalgebra
of $\mathfrak{so}(25)$, with $e_1(h)=2$ and $e_i(h)=0$
for $i>1$. Let $\mathfrak t'=\ker e_{12}$.

\begin{lemma}\label{lem:restriction}
Suppose two finite-dimensional $\mathfrak{so}(25)$ modules
$V,V'$ have equal characters on $\mathfrak t'$. If every
highest weight $\nu$ of $V'$ has $\nu_{11}=0$, then
$V\cong V'$.
\end{lemma}
\begin{proof}
Use lexicographic order. Every weight $\alpha$ of
$\VV_\lambda$ satisfies $\alpha\le\alpha^+\le\lambda$,
where $\alpha^+$ is obtained by sorting its absolute
coordinates. The second inequality follows from the positive
root order. The largest restricted highest weights of $V$
and $V'$ therefore have identical first eleven coordinates.
If they are $\lambda$ and $\nu$, the assumption gives
$\lambda_{11}=\nu_{11}=0$, and dominance forces
$\lambda_{12}=\nu_{12}=0$; hence $\lambda=\nu$.

No other weight can contribute to this largest restricted
weight. For if $\alpha$ has those first eleven coordinates,
then $\alpha\le\alpha^+\le\lambda'\le\nu$ for its
constituent highest weight $\lambda'$. The first eleven
coordinates of $\alpha^+$ must agree with $\nu$. Its
eleventh and twelfth coordinates are zero, so
$\alpha^+=\lambda'=\nu$ and $\alpha_{12}=0$.
The multiplicity of the restricted weight is therefore exactly
the multiplicity of $\VV_\nu$. Subtract this common summand
and continue by induction.
\end{proof}

The hypothesis is necessary: the sums of
$(25-2i)\bigwedge^i\CC^{25}$ over even and odd $i$, for
$0\le i\le12$, have the same restricted character. Nor do
the degree and Hodge gradings alone distinguish representations:
$\VV_{(1,1,1)}\oplus230\VV_{(1)}$ and
$21\VV_{(1,1)}\oplus1750\VV_{(0)}$ agree on that
two-dimensional torus. These examples also test the scope
of the character argument.

\subsection{The Hodge-theoretic hypothesis}
Let $(S',\beta)$ be a Picard-rank-one K3 surface of degree four,
fix $\chi_0\ne0$, and put
\[
 w=(0,\beta,\chi_0),\quad v'=(0,2\beta,2\chi_0+1),\qquad
 M'=M_\beta(2w),\quad N'=M_\beta(w),\quad X'=M_\beta(v').
\]
The last two spaces have types $K3^{[3]}$ and $K3^{[9]}$,
and $M'$ is locally trivially deformation equivalent to $M$
\cite[Theorem~1.7]{pr23}. The additional hypothesis is:
\begin{enumerate}
\item[(H3)] For every such $(S',\beta)$ there is an isomorphism
of graded rational Hodge structures, with usual degrees,
\begin{equation}\label{eq:H3}
 H^\bullet(X')\cong\IH^\bullet(M')\oplus
                 \bigwedge^2H^\bullet(N')[-6](-3).
\end{equation}
Here $[-6]$ raises degrees by six, and $(-3)$ is the Tate twist.
\end{enumerate}
We do not prove (H3). Its proposed source is the mixed-Hodge-module
form of the BPS theorem
\cite[Theorems~9.2 and~10.1, Corollary~3.6,
Propositions~6.10 and~B.14]{dhsm22} and the independence theorem
\cite[Theorem~5.1]{dhksv26}. The first splits the BPS object
of $2w$ into the intersection object and the antisymmetric
direct-sum contribution from $N'\times N'$. The second
identifies the perverse direct images under the support maps
for different Euler characteristics as mixed Hodge modules.
Saito's decomposition theorem \cite[(4.5.4)]{saito90} then
gives \eqref{eq:H3}; weights impose the displayed twist.
This deduction uses the mixed-Hodge-module assertion of the
second theorem, whose maps are constructed by Hecke operators.
Equality of Hodge numbers, as supplied by \cite{gwz26}, is
insufficient for what follows.

\begin{theorem}\label{thm:conditional}
Assume (H3). For the eighteen-dimensional moduli space in
\eqref{eq:family}, $\IH^\bullet(M;\CC)\cong\mathcal R$
as a module over its LLV algebra.
\end{theorem}
\begin{proof}
Choose a smooth algebraic family of degree-four polarised K3
surfaces with arithmetic monodromy on $\beta^\perp$
\cite[Corollary~6.4.7]{huybrechts16}. Form the three relative
moduli spaces with vectors $2w,w,v'$. Restrict to the open
set where the polarisation is \emph{general with respect to}
these vectors in the sense of \cite[Definition~2.8]{pr23};
the required openness is \cite[Proposition~2.14]{pr23}.
This is the specific condition used here, and is not an
identification with every wall-based genericity convention.
The open set contains the Picard-rank-one fibres.

Shrink further to a smooth connected dense open base $T^\circ$
where the singular family is topologically locally trivial
with its Whitney strata \cite[Corollaire~(5.1)]{verdier76}
and the two primitive families are smooth. Intersection
cohomology then forms a local system with its integral
structure, pairing and ordinary-cohomology action. In each
degree it underlies a polarizable pure variation after a
further shrinking. To justify that its fibre structure is the
usual intersection Hodge structure, apply Saito's projective
direct-image theorem and generic variation theorem
\cite[Th\'eor\`emes~5.3.1 and~5.4.3]{saito88} to the
intersection Hodge module of the total space. Base change
\cite[(4.4.3)]{saito90} identifies its restriction to a fibre;
local triviality identifies the underlying perverse sheaf
with the intersection complex. Its restriction to the regular
locus is the constant variation, which determines the
intersection extension and its Hodge structure, with the
degree shifts and possible Tate normalisation.

Choose a Hodge-generic point of this base whose surface has
Picard rank one. Such points exist outside a countable union
of proper analytic subsets. Let $\Gamma=\pi_1(T^\circ)$.
Removing proper complex subvarieties from a smooth base
surjects on its original fundamental group, so the image of
$\Gamma$ on $\beta^\perp$ remains arithmetic and Zariski
dense in $\SO(\beta^\perp_{\CC})$.

The Mukai isometries for $N'$ and $X'$ identify their second
cohomologies with $w^\perp$ and $(v')^\perp$
\cite[Theorem~1.7]{markman08}. These identifications are
locally constant \cite{ogrady97,markman08,opr24}. Their common
subspace $\beta^\perp$ has dimension $21$, and monodromy is
trivial on its two-dimensional complements, consisting of
vectors $(r,a\beta,s)$. The diagonal restriction
$H^2(M')\to H^2(N')$ is a similitude by the argument in
Proposition \ref{prop:og10}. This assertion may be checked
using the full locally trivial monodromy of $M'$; that group
also transports its intrinsic diagonal stratum and acts
isometrically on $H^2(N')$. Restriction is a morphism of local
systems in our chosen family. It identifies a corresponding
nondegenerate $21$-dimensional subspace $W\subset H^2(M')$
with $\beta^\perp_{\CC}$ and identifies their $\Gamma$ actions.

Transport Proposition \ref{prop:spin} from $M$ to $M'$;
apply its smooth version to $N'$ and $X'$. A finite-index
subgroup of $\Gamma$ acts on their cohomologies through
the corresponding spin images. All three actions project
to the same action on $\beta^\perp$. Their fibre product
over $\SO(\beta^\perp_{\CC})$ is a finite cover of that
group; its identity component is the diagonal image of
$\operatorname{Spin}(\beta^\perp_{\CC})$. Density of
the projection implies that this diagonal image is the
connected monodromy group on
\[
 \mathcal V=\IH^\bullet(M')\oplus
           \bigwedge^2H^\bullet(N')\oplus H^\bullet(X').
\]

Degree by degree these are polarizable pure variations; their
direct sum is a graded-polarizable variation with several
weights. Andr\'e's monodromy theorem
\cite[Lemma~4, Theorem~1]{andre92} puts its connected monodromy
group in the derived group of the generic Mumford--Tate group.
It therefore fixes every rational tensor of type $(p,p)$:
after a Tate twist such a tensor is Hodge, and the derived
group acts trivially on Tate structures. The components of
the isomorphism \eqref{eq:H3}, extended by zero to operators
on the indicated direct sum, are such tensors. Hence (H3)
is equivariant for this diagonal spin group at the chosen
generic point. No isomorphism of variations is assumed.

Include $h$ with a Cartan subalgebra of $\mathfrak{so}(W)$.
The resulting subalgebra is $\mathfrak t'=\ker e_{12}$.
Extend the identifications on $W$ to complex similitudes of
the full second cohomologies; the complementary quadratic
spaces have dimension two, so such extensions exist.
After centring the gradings, \eqref{eq:H3} now says that
the restricted character of $\IH(M')$ is that of
$H(K3^{[9]})-\bigwedge^2H(K3^{[3]})$, hence that of
$\mathcal R$. Every highest weight in Table
\ref{tab:representation} has at most three nonzero entries.
Lemma \ref{lem:restriction} identifies the full modules.
Parallel transport from $M'$ to $M$ completes the proof.
\end{proof}

This last argument uses (H3) at a Hodge-generic fibre only.
It does not identify Saito's Hodge decomposition with the
type decomposition of harmonic forms for the metric on $M$.
